\scanpage{23}
\mspage{\unsure{6}}
\uline{Thm} Considérons les \inserted{deux} catégories \deleted{\unsure{définies ainsi}} :
a) celle des \deleted{triples}\inserted{systèmes}
$(\widetilde X^*,\dot X,\mathcal E,Y,f)$, où
$\widetilde X^*$ espace top., $\dot X$ partie fermée, $\mathcal E$ un germe
de voisinage cylindrique de $\dot X$ dans $\widetilde X^*$,
$Y$ esp. top., $f:\dot X\longrightarrow Y$ application continue ;
b) celle des systèmes $(X,Y,\mathcal T)$, où
$X$ esp. top., $Y$ partie fermée, $\mathcal T$ un germe
de voisinage cônique de $Y$ dans $X$. \deleted{\unsure{Alors} Alors}
\deleted{\unsure{le}} Alors le foncteur naturel \inserted{ci-dessus} de a)
dans b) est une équivalence de catégories.

\margin{NB Il y a un énoncé \unsure{analogue} en termes de structures de
voisinages tubulaires \uline{stricts} (au lieu des germes
de voisinages tub.) \unsure{à} \uuline{voir !}
Les catégories sont \unsure{$\simeq$}.
Les morphismes ici sont les isos, pour simplifier !}

Rappelons que
\[
X=\widetilde X^*\amalg_{\dot X}Y
\]
\margin{\uline{NB} $(\widetilde X^*,\dot X)$ joue le rôle d'un
``découpage cônique'' de $X$ suivant $(Y,\mathcal T)$ \unsure{ci-dessus}.}

\uline{NB} Une partie \inserted{$X'$} de $X$ revient donc à une
partie $\widetilde X'^*$ de $X^*$, $Y'$ de $Y$, ayant $=$
images inverses dans $\dot X$, et $X'$ est fermé ssi
$\widetilde X'^*$ et $Y'$ le sont. \deleted{Pour que $X'$ soit}
\inserted{Pour que $X'$ soit}
compatibles avec le germe de structure
cônique $\mathcal T$, \deleted{il \unsure{f.} \inserted{une description} que $X'_{\unsure{*}}$}
\inserted{il f. et s. que l'on ait}
a) $\widetilde X'$ compatible avec le germe de structure
cylindrique $\mathcal E$ de $(\widetilde X,\dot X)$.
Il y a lieu alors de remplacer $\widetilde X'$ par
$\inserted{\widetilde X'_1}\,\overline{X'\cap X^*}$
(où $X^*=\widetilde X\setminus\dot X$), et le germe de $\widetilde X'_1$ au vois.
\margin{22}
